\scanpage{122}
\deleted{\unsure{Voir|Voici|Quid} P} avec deux sous-alg. \deleted{\unsure{H}}étales (\deleted{diago} scindées) $\overline H,\overline H'$ correspondant à l'inv. $\rho_0$. Quitte à faire encore ext. biquadratique, \deleted{$H\unsure{\simeq|=}\overline H,H'$} $\overline H\simeq H$, $\overline H'\simeq H'$, ou encore \deleted{\unsure{Donc}} \deleted{\unsure{Cela est}}\inserted{OPS $H\unsure{,|\ne}H'\subset A$.} Soit alors
\[
\tau=\operatorname{Tr}uv,
\]
si on avait $t=\tau$, on aurait gagné. On sait que
\[
\deleted{b^2-4\unsure{c|\cdot}}\frac{\delta\delta'}{Q(t)}=\frac{\delta\delta'}{Q(\tau)}\qquad(=\rho_0)
\]
donc
\[
Q(t)=Q(\tau)\deleted{\unsure{=}}\qquad(=\sigma)
\]
\deleted{\unsure{Or les|On les}} ou encore, posons \deleted{Supposons $k$ corps alg. clos. Comme}
\[
\begin{aligned}
\deleted{\unsure{t^2}}\quad R(T)&=T^2-bb'T+r\\ \text{où}\qquad r&=(b^2c'+b'^2c-4cc'-\sigma)
\end{aligned}
\]
$t$ et $\tau$ sont des solutions de l'équation $R(T)=0$. Si donc on avait $\tau\ne t$, on aurait
\[
t+\tau=\deleted{\unsure{+}}bb'\qquad\text{donc}\quad\tau=bb'-t
\]
Or considérons
\[
u'=\operatorname{Tr}u\cdot1-u=b-u
\]
% archive sheet number: 121
